% Manuscript-ready seed-level ablation table for the frozen LAM-JEPA ARC study.
% Derived only from retained five-seed accuracies on the fixed 295-row validation cohort.
% Does not authorize any claim beyond the bounded frozen-protocol result.

\begin{table}[t]
\centering
\caption{Seed-level correct-count decomposition for the frozen mechanism ablations. Each seed uses the same 295-row validation cohort. The reviewed retained runs are constant classifiers, so sparse count differences should not be interpreted as stable input-conditional mechanism effects.}
\label{tab:seed-level-ablation-counts}
\begin{tabular}{rrrrrrr}
\toprule
Seed & Full & No planner & $\Delta_{\mathrm{planner}}$ & No target & $\Delta_{\mathrm{target}}$ & Shuffled \\
\midrule
1 & 71 & 71 & 0 & 71 & 0 & 78 \\
2 & 78 & 78 & 0 & 78 & 0 & 71 \\
3 & 78 & 71 & +7 & 83 & -5 & 83 \\
4 & 71 & 71 & 0 & 71 & 0 & 78 \\
5 & 78 & 78 & 0 & 83 & -5 & 78 \\
\bottomrule
\end{tabular}
\end{table}

% Suggested accompanying prose:
% Seed-level integer outcomes make the five-seed limitation explicit. Over the
% fixed 295-row validation set, full-minus-no-planner correct-count differences
% were [0, 0, +7, 0, 0], while full-minus-no-target differences were
% [0, 0, -5, 0, -5]. Thus the planner aggregate is driven by a single seed and
% the target-path aggregate by two seeds. Because the externally reviewed
% retained runs are constant classifiers, these sparse count shifts are
% consistent with changes in the class selected under collapse rather than
% evidence of stable input-conditional planner or target-path benefit.
